\section{Conditional labels and entropy}\label{sec:trees}

Lemma~\ref{lem:hierarchy} counted the total pairing volume supplied by
successive labels.  We now keep track of which pair of sides supplies
each contribution.  The input is a finite tree describing which label
each pair can read; the output is a matrix-multiplication bound in
terms of the three accumulated conditional entropies.  Preserving
previously read labels on every side is essential when the next
reader pair changes.

\begin{definition}\label{def:label-tree}
A \emph{readable label tree} for the support of a tensor is a finite
rooted partition tree of the support whose leaves are in bijection
with the supported monomials.  Each monomial follows exactly one path.
At each nonterminal node, conditional on the path to that node, the
next child is a function of either of two specified side coordinates
separately.
The reader pair may vary with the node.
\end{definition}

For a probability law \(\pi\) on the leaves, write \(\Prb(\nu)\) for
the mass passing through a node \(\nu\).
Assign to each pair \(UV\) the rate
\[
 r_{UV}(\pi)=
 \sum_{\nu:\,\text{reader pair }UV}
 \Prb(\nu)\,\HH(\text{child}\mid\nu),
 \qquad
 \HH(p_1,\ldots,p_q)=-\sum_i p_i\log p_i.
\]
Zero-mass contributions are zero.
The sum of the three rates is \(\HH(\pi)\).

This identity is the entropy chain rule \cite[Section~6]{shannon1948}.
The exact conditional counts below use the method of types
\cite[Section~II]{csiszar1998}; the tensor restrictions and their full
auxiliary costs are proved explicitly.

\begin{theorem}[Conditional-label entropy bound]\label{thm:entropy}
If a tensor of cost \(e^C\) has a readable label tree, then
\begin{equation}\label{eq:entropy-bound}
 W(r_{XY}(\pi),r_{XZ}(\pi),r_{YZ}(\pi))\le C
\end{equation}
for every leaf law \(\pi\).
\end{theorem}

\begin{proof}
First suppose \(\pi\) is rational and take \(n\) with all \(n\pi_s\)
integral.  Prescribe the resulting counts at every node.
Process nonterminal nodes in any order in which a parent precedes
its children.

The induction invariant is that the current tensor is a direct sum
over all processed word choices, and that each side identifies the
entire processed prefix.  Within a summand, the tensor is the original
word-support fiber times the dot factors already introduced.
At the root this invariant is trivial.

Consider a node \(\nu\) occurring in \(s\) positions, with child counts
\(s_1,\ldots,s_q\).
Within a prefix summand, all sides know these positions.
The two readers can therefore project to child words with those counts
and apply \cref{lem:separator} with
\[
 K_\nu=\binom{s}{s_1,\ldots,s_q}.
\]
The new tag makes the chosen child word visible on the third side.
All previous dot indices are carried unchanged, so the invariant
continues.

Exactly one auxiliary tensor is used for this node, even if its
positions vary between prefix summands.
Indeed,
\[
 \left(\bigoplus_p T_p\right)\otimes A_{K_\nu}
   =\bigoplus_p(T_p\otimes A_{K_\nu}).
\]
The prefix blocks are disjoint on all three sides, so the restrictions
within them may depend on \(p\).
The number \(K_\nu\) depends only on the prescribed counts, not on
the positions of a particular prefix.
If \(s=0\), take \(K_\nu=1\) and do nothing.

At the end, every surviving original word is a specified monomial
word with leaf histogram \(n\pi\).
Its scalar coefficient is nonzero and can be absorbed by a local
rescaling in that summand.
The remaining tensor is a matrix product of dot factors, with pair sizes
\[
 l_{UV}(n)=\prod_{\nu:\,\text{reader pair }UV}K_\nu.
\]
These sizes do not depend on the particular surviving word.
The number of summands is
\begin{equation}\label{eq:multinomial-telescope}
 M(n)=\frac{n!}{\prod_s(n\pi_s)!}=\prod_\nu K_\nu.
\end{equation}
The second equality is the cancellation of child factorials with
parent factorials; it also counts the successive choices of words.
Every such word survives all the restrictions.

The full cost of this construction is
\[
 e^{nC}\prod_\nu |G_{K_\nu}|.
\]
For the fixed finite tree and law, a node with positive conditional
entropy has \(\log K_\nu\) linear in \(n\) up to \(O(\log n)\);
a deterministic node has \(K_\nu=1\).
The logarithmic factorial estimate and
\cref{lem:separator} therefore give
\[
 \frac{\log l_{UV}(n)}n\longrightarrow r_{UV}(\pi),\qquad
 \log\!\left(\prod_\nu|G_{K_\nu}|\right)-\log M(n)=o(n).
\]
Apply \cref{lem:rectangular-sum} at each fixed \(n\), divide by \(n\),
and use homogeneity and continuity from \cref{lem:rank-exponent}.
This proves \eqref{eq:entropy-bound}.
Finally approximate an arbitrary law by rational laws.
All the finite entropies, and \(W\), are continuous.
\end{proof}

\begin{corollary}[Balancing two axes]\label{cor:balance}
Fix a side \(L\).
Let \(S\) be the sum of the two rates incident to \(L\), and let \(T\)
be the remaining rate.  Then
\begin{equation}\label{eq:balanced-rate}
 W(S/2,T,S/2)\le C.
\end{equation}
For \(C>0\), the normalized pairs \((S/C,T/C)\) from different constructions
may be combined convexly and decreased coordinatewise while retaining
the normalized bound.
\end{corollary}

\begin{proof}
Average \eqref{eq:entropy-bound} with the estimate exchanging the two
incident axes, using permutation invariance, subadditivity and
homogeneity.  The same properties justify convex combinations after
normalizing each bound by its cost.  Monotonicity permits coordinatewise
decrease, and continuity permits closure.
\end{proof}

The auxiliary cost in \cref{thm:entropy} is not suppressed.
Its exponential contribution is precisely what the multiplicity
\eqref{eq:multinomial-telescope} cancels after
\cref{lem:rectangular-sum}.
This exact accounting is what permits many successive conditional
separations.
